% DRAFT (Claude, 2026-11-11, Anders): the fourth paragraph of Anders' own 5.2 ("Scheduling Hierarchy and Valuation
% Perspective"), rewritten to say more about each study. Only this paragraph; the rest of Anders' text is unchanged.
% Model (precedent rule): bore-brensholm-2025 P23 (GROUP: the shared thing with a group citation; each study with aim,
% method, markets and horizon in one or two sentences) with the closer of bore-2026 P33 (a limitation of the group).
% Third version after two independent reviews: opener limited to these studies and "expected" dropped (Ak averages
% realised years); year counts removed (no numbers before the caveat); Ak's gain hedged; Helseth 2016 given its reason
% (what an added revenue source does to the water values) and its short-term consequence left to the next paragraph.
% Every citation's passage in CHANGES-5-2.md (2026-11-11 section).

Medium-term scheduling studies such as these take prices as input and
maximise the profit of a single producer
\citep{ak2019quantifying,alic2023optimal,helseth2016}.
% src: ak2019quantifying p. 741 (price-taker; realised hourly prices; average annual revenue); alic2023optimal p. 17; helseth2016 p. 934
\citet{ak2019quantifying} quantify the gain from running an existing
Turkish cascade as pumped storage, using monthly rule curves and the hourly
prices of past years.
% src: ak2019quantifying p. 741 (monthly rule curves with hourly prices within each month; past realised prices), p. 744 (cascade), p. 748 (2013-2017)
In their case the monthly use of water hardly changes with pumping. The
gain may therefore come from the hours of the day in which the plant
turbines and pumps, net of the pumping cost.
% src: ak2019quantifying p. 751 ("similar amounts of water are turbined monthly under both systems, however PSH alters the hours of the day ... may result in significant gain"); Table 2 p. 747 (net revenue)
\citet{alic2023optimal} also add a hypothetical pump to the upper station
of a Norwegian cascade. Their stochastic dynamic programme has weekly
stages, a joint price--inflow Markov chain and environmental restrictions.
% src: alic2023optimal p. 27 ("cannot currently operate as a PSHP"), p. 29 (38 MW), pp. 11-12 (SDP, weekly stages, Markov chain of inflow and price), pp. 16-22 (MEF, SCWA, ramping)
The pump raised their revenue even under the ramping limits that lowered
the revenue of the conventional plant.
% src: alic2023optimal p. 32 (Case A +2.35 %, Case B -1.01 %, Case C +0.47 %; "overcome the limitations imposed by the ramping constraints")
\citet{helseth2016}, in contrast, consider a conventional watercourse
whose price-taking owner sells energy and reserve capacity, which shows
what an added source of income does to the water values.
% src: helseth2016 p. 934 (price-taking, energy and capacity), p. 939 (weekly stages); p. 941 (Fig. 7, water values with and without reserves)
In their case a significant part of the capacity income came without a
changed generation schedule, but the water values changed.
% src: helseth2016 p. 940 ("A significant part of the profit from capacity sales ... without changing the generation schedule"); p. 941 ("water values generally seems to be lower")
Although the studies differ in configuration and markets, each allocates
the water by the market income it can earn, and none of the three varies
the pumping capacity.
% src: ak2019quantifying Table 1 p. 746 (fixed reversible-turbine capacities); alic2023optimal p. 29 (one 38 MW pump); helseth2016 (no pump)
